%% 第六章
\chapter{能力优化方法与实验}
\chaptermark{能力优化方法与实验}

\section{UAV-DualCog-Train训练集构建}

第五章的评测结果表明，当前多模态大模型在 UAV-DualCog 上的主要短板并不是完全无法回答无人机场景问题，而是难以同时保持语义答案、显式空间证据、跨媒介稳定性以及自我状态认知与环境情况认知之间的平衡。围绕这些问题，本文进一步实现了面向能力优化的训练数据与训练脚本，并将其组织为 UAV-DualCog-Train-V1。

训练集使用当前未纳入测试基准的 6 个场景构建，分别为 \texttt{env\_1}、\texttt{env\_3}、\texttt{env\_5}、\texttt{env\_15}、\texttt{env\_18} 和 \texttt{env\_22}。这些场景与第五章评测所使用的测试场景在场景层面严格隔离，从而降低场景泄漏造成的泛化偏差。训练集沿用 UAV-DualCog 的问题定义、输入格式、输出字段和评价口径，但当前能力优化阶段优先聚焦四类图像任务，即地标相对位置推理、未来观察预测、自我参照位置推理和地标驱动动作决策。这样设计有两个直接原因：一方面，第五章已经表明图像侧普遍存在“语义相对较强、空间证据明显偏弱”的现象；另一方面，图像任务天然适合把“选项判断+目标框证据”写成统一的结构化监督目标。

训练数据由 \texttt{prepare\_train\_dataset.py} 自动从各场景的 Stage 4 最新清单文件中导出，主要脚本入口见\autoref{app:optimization-entry}。每条训练样本包含样本标识、任务类型、难度、候选选项、标准答案、标准化目标框以及参考图和查询图等关键信息；随后，脚本把这些字段统一转换为训练记录，包括文本提示 \texttt{prompt}、监督目标 \texttt{response} 以及 0 到 1000 整数域上的标准化目标框 \texttt{answer\_bbox\_xyxy\_1000}。其中监督目标固定写为两行：
\begin{verbatim}
Answer: <option>
Bounding Box: [x1, y1, x2, y2]
\end{verbatim}
这一设计与第五章的结构化评测协议保持一致，也便于后续在 SFT 和 GRPO 两个阶段共享同一套输出解析逻辑。

在数据完整性方面，脚本会自动检查每个场景的最新清单是否存在、样本数量是否为正、四类任务与两档难度是否全部覆盖，以及图像路径、目标框和选项字段是否存在异常。根据完整性核查报告的结果，当前 UAV-DualCog-Train-V1 共包含 11232 条样本，其中训练集 10104 条，验证集 1128 条，图像路径缺失、目标框异常和选项字段缺失均为 0。训练集与验证集采用按任务类型和难度分桶后的稳定轮转划分策略，即每 10 条样本取 1 条进入验证集，其余进入训练集。这样既避免了完全随机划分带来的波动，也能在不同任务键上保持近似一致的 9:1 比例。

\section{监督微调方法实现}

能力优化的第一阶段采用监督微调。该阶段的目标不是改变基座模型结构，而是让模型先稳定学习 UAV-DualCog 图像任务中的结构化输出协议，并建立“答案与证据同时生成”的基础行为。当前实现选择 \texttt{Qwen/Qwen3.5-4B} 作为基座模型，一方面是因为 Qwen 系列在第五章评测中表现具有代表性，另一方面是因为 4B 规模便于在现有计算资源下完成多轮实验和多阶段对比。

SFT 由 \texttt{run\_sft\_lora.py} 实现。脚本首先读取上一步生成的训练与验证 JSONL 文件，然后将每条样本拼接为单段因果语言模型训练文本，即“任务提示词 + 结构化监督答案”。在提示词构造上，脚本显式写入任务类型、问题文本、候选选项、参考图路径和查询图路径，并要求模型“严格返回两行”。当前公开训练脚本并未在该阶段直接引入像素级视觉编码器，而是以结构化文本样本完成训练链路验证和方法实现。这一点与论文中方法定位保持一致，即本章首先解决结构化证据输出与后训练闭环问题。

为兼顾资源约束与可复现性，SFT 阶段采用 LoRA 参数高效微调。脚本中默认超参数为：学习率 $5\times10^{-5}$、训练轮数 3、单卡批次大小 1、梯度累积步数 32、最大序列长度 2048、LoRA 秩 16、LoRA 系数 32、LoRA dropout 0.05。LoRA 目标层设置为 \texttt{q\_proj}、\texttt{k\_proj}、\texttt{v\_proj}、\texttt{o\_proj}、\texttt{gate\_proj}、\texttt{up\_proj} 和 \texttt{down\_proj}，使用 \texttt{Transformers}、\texttt{PEFT} 和 \texttt{Trainer} 完成训练。训练结束后，脚本会保存 LoRA 适配器、训练指标文件，并可选导出 \texttt{merged\_model} 作为下一阶段 GRPO 的初始化模型。

从训练目标上看，SFT 阶段主要解决三个问题。第一，模型能否稳定输出可解析的两行结构，而不出现字段缺失、答案格式漂移或目标框无法解析的情况。第二，模型能否在四类图像任务上建立基本任务适应能力，使语义判断与候选选项约束对齐。第三，模型能否为下一阶段基于奖励的细粒度优化提供一个足够稳定的初始化点。与单纯提升答案准确率相比，这一阶段更强调训练接口、监督目标和输出协议的一致性。

\section{轻量GRPO优化方法实现}

第二阶段采用轻量组相对策略优化（Group Relative Policy Optimization, GRPO）。与直接从基座模型开始强化训练不同，本文实现中始终以上一阶段的 SFT 结果作为初始化。这样做的目的，是把有限训练预算集中用于修正“答对但框错”“格式正确但证据不可信”这类更贴近 UAV-DualCog 核心问题的失败模式，而不是反复纠正最基础的输出格式。

GRPO 由 \texttt{run\_grpo\_light.py} 实现，核心奖励函数由 \texttt{rewards.py} 中的 \texttt{total\_reward} 计算。当前实现采用三类奖励项的加权求和：
\begin{equation}
R = w_f R_{\text{format}} + w_s R_{\text{mcq}} + w_b R_{\text{bbox}}.
\end{equation}
其中，$R_{\text{format}}$ 用于判断模型输出是否能够被解析为规定的两行结构，$R_{\text{mcq}}$ 用于衡量多项选择答案是否正确，$R_{\text{bbox}}$ 用于衡量预测目标框与真实目标框之间的一致性。当前实现中，目标框奖励以标准化整数域上的框坐标为基础，通过几何重合质量映射为空间证据得分。

为了降低训练震荡和奖励作弊风险，GRPO 采用两阶段渐进式设计。Phase A 对应 \texttt{semantics\_first}，只启用格式奖励和语义奖励，默认权重为 $(w_f,w_s,w_b)=(1.0,1.0,0.0)$；Phase B 对应 \texttt{with\_bbox}，在前一阶段基础上加入空间证据奖励，默认权重为 $(1.0,1.0,1.0)$。这一设计与第五章结论形成直接对应：先稳定“答得对、写得稳”，再进一步约束“证据也要对”。由于训练脚本支持从适配器检查点继续启动，若输入的是仅含 LoRA 适配器的检查点，脚本会自动完成一次合并，生成临时的全量模型目录，再交给 \texttt{GRPOTrainer} 使用，从而减少多阶段衔接中的人工处理成本。

当前实现的推荐初值同样以资源受限条件为前提。Phase A 默认学习率为 $2\times10^{-6}$，Phase B 默认学习率为 $1\times10^{-6}$；两个阶段均采用生成组大小 2、单卡批次大小 1、梯度累积步数 16、最大提示长度 1536、最大生成长度 128。脚本使用 \texttt{TRL} 中的 \texttt{GRPOTrainer} 完成训练，并将奖励函数直接作用在每次生成结果上。相比全尺寸强化学习配置，这是一条更轻量、可落地、与当前硬件条件相匹配的后训练路径。

\section{训练流程、工程验证与当前状态}

为保证训练过程可复现、可追踪，项目进一步实现了全流程启动脚本 \texttt{run\_full\_pipeline.sh}。该脚本统一管理数据输入、环境检查、单卡或多卡启动方式、日志目录、检查点目录和输入快照目录。若训练开始前检测不到训练集或验证集 JSONL 文件，脚本会自动调用 \texttt{prepare\_train\_dataset.py} 重新生成；若文件已存在，则会把当前使用的训练集、验证集以及全量样本快照复制到本次运行目录中，便于后续复现实验来源。完成快照后，脚本依次调用 SFT、GRPO Phase A 和 GRPO Phase B，并在每个阶段结束后检查 \texttt{merged\_model} 是否生成，从而保证后续阶段不会建立在缺失或不完整的中间产物之上。

在正式全量训练之前，本文已经完成了三类工程验证。第一类是链路级 smoke 验证，包括训练集准备、SFT dry-run、GRPO dry-run 以及迁移评测命令链路检查。第二类是功能级验证，包括一次非 dry-run 的 1 步 CPU GRPO 小样本训练，以及一次 2 样本 SFT 前向可用性验证，验证结果表明训练脚本、损失计算和奖励接口均可正常工作。第三类是端到端 Tiny GPU Run 验证：在 64 条训练样本、最多 20 步的设置下，SFT、GRPO Phase A 和 GRPO Phase B 三段流程均已完整跑通，SFT loss 与验证 loss 呈稳定下降趋势，GRPO Phase A 的平均奖励从 1.312 上升到 1.425，Phase B 的目标框奖励逻辑也已正常触发并完成产物落盘。需要说明的是，这些结果仅用于证明训练链路和方法实现是可运行、可衔接、可收敛的，不作为正式能力优化实验结果写入后文分析。

截至 2026 年 5 月 19 日，UAV-DualCog 的能力优化部分已经完成训练数据导出、完整性核查、SFT 实现、两阶段轻量 GRPO 实现、全流程调度脚本编写以及 smoke 和 tiny 规模验证，正式全量训练正在进行中。因此，本章的方法与实验设置部分已经从“拟开展方案”转变为“已实现并已验证链路的训练方法”，而正式对比实验结果、消融结果和训练后案例分析将在训练结束后补充。

\section{实验设置与结果分析预留}

能力优化实验仍将保留四组核心对比模型，即原始 Qwen 3.5 4B 基座模型、SFT 模型、SFT+轻量 GRPO（格式+语义奖励）模型和 SFT+轻量 GRPO（格式+语义+空间奖励）模型。图像任务将继续报告选项准确率、BBox Acc@50IoU 和 mIoU，并分别统计自我状态认知与环境情况认知两个分支；视频任务将继续使用语义识别、计数准确率、时间区间 F1@0.5 和 mTIoU 等指标，用于检验图像侧证据感知训练能否向视频侧迁移。

由于正式全量训练尚未结束，本节暂不写入最终定量结果和定性案例。训练完成后，本章将重点回答以下问题：SFT 是否已经显著提升四类图像任务上的结构化输出稳定性与基础准确率；轻量 GRPO 是否能在 SFT 基础上继续提升语义与证据一致性；加入目标框奖励后，改进主要体现在哪些任务、哪些认知分支和哪些指标上；自我状态认知与环境情况认知之间的差距是否缩小；图像任务上的证据感知训练是否能够迁移到视频任务。

\section{本章小结}

本章根据第五章的评测发现，系统整理了 UAV-DualCog 能力优化部分的真实实现状态。首先，本文完成了 UAV-DualCog-Train-V1 的数据导出、结构化样本组织、完整性核查和稳定划分。其次，本文完成了基于 Qwen 3.5 4B 的 LoRA 监督微调实现，以及基于格式奖励、语义奖励和空间证据奖励的两阶段轻量 GRPO 实现。再次，本文通过 smoke 验证、功能验证和 tiny GPU 端到端验证，确认了数据、脚本、奖励函数和阶段衔接链路的可用性。由此，第六章已经具备完整的方法论和工程实现基础；待正式训练完成后，仅需补充最终结果表、案例分析和消融实验，即可与第五章的能力评测结论形成完整闭环。
